\section{The bypass in a switch pipeline (H4)}
\label{app:switch}

This appendix spells out what a switch must do to implement the bypass, and what that does and does not say about its cost. No hardware pipeline was compiled or measured, and the simulator adds no processing delay for route selection, so our results measure the effect of the changed paths, not the cost of choosing them.

\paragraph{Required decisions.} These follow from the simulator's routing (\file{src/tcp-bypass/bypass-routing.cc}), with local delivery handled first:

\begin{center}
\small
\begin{tabularx}{\linewidth}{@{}l L@{}}
\toprule
Packet & Decision \\
\midrule
Destination host is on this router & Deliver locally. \\
Marked (small flow), at any router & Pick one of the neighbors closer to the destination, by flow hash. \\
Unmarked & Spraypoint's spray or point decision, unchanged. \\
\bottomrule
\end{tabularx}
\end{center}

The inputs are the destination prefix, the ingress port class, and the small-flow mark: DSCP AF11, read as \texttt{(ToS \& 0xfc) == 0x28} so the ECN bits are kept. A target can carry the classification as one metadata bit after a DSCP match. A consistent flow hash keeps each flow on one path. The design assumes the sending host knows a flow's size and marks every packet of it, including TCP control packets; it does not ask the switch to count bytes per flow or predict flow sizes, which would be a different design with different state and reordering costs.

\paragraph{A logical implementation.} RNG already selects its table by ingress port: server ingress sprays, fabric ingress points, and local destinations are delivered (§5 of~\cite{rng}). The bypass adds a classification and one optional lookup before that baseline:
\begin{verbatim}
if destination is local:      deliver locally
else if small:                use the shortest-path ECMP group for the destination
else:                         apply the baseline spray or point route
\end{verbatim}
This is a forwarding specification, not compiled P4 or a count of physical stages. On a device that can fold the classification into an existing table or use spare resources in its normal traversal, the extra decision need not cost another pass. A dependent lookup may instead need an extra physical stage, and if resources run out the program may not fit or may need recirculation. Table widths, action capacity and existing occupancy all have to be checked on a real target.

\paragraph{State.} One prefix identifies each router's hosts, so the bypass needs at most $n - 1$ extra destination entries per router (255 in our simulation, 999 at the paper's scale), each pointing to an ECMP group of the neighbors closer to that destination. A group can hold up to $d$ members, so membership is bounded by $O(nd)$ per router, not by Spraypoint's $h = 2$; shared groups can reduce it. These groups compete with Spraypoint's pointing groups for ECMP memory (§7.4 of~\cite{rng}). The adjacent bypass of \secref{sec:bypass} needs only $d$ entries. The simulator's all-pairs breadth-first search (\file{src/tcp-bypass/shortest-paths.cc}) is how routes are computed, not state each switch stores or work it does per packet.

\paragraph{What follows about latency.} If a target adds at most $\delta$ processing delay per switch, a packet crossing $k$ switches pays at most $k\delta$, under the resource assumptions above. For scale, every fabric hop the bypass removes saves 500\,ns of propagation in our simulation, before serialization and queueing. So H4 holds conditionally: the bypass is a destination-based override, not per-flow state, and a normal pipeline mapping would make its cost small next to the delay it saves. The proposal's stronger claim, exactly one extra lookup with negligible cost on every commodity switch, remains unverified. Verifying it would take implementing the baseline and the bypass on a named switch architecture, checking the compiler's resource report at the intended scale, testing the forwarding behavior with packets, and measuring forwarding latency and rate with and without the bypass.
